% TOP_PROBLEM: {"id":30007563,"problem_number":"LOCAL-30007563","title":"Learning-based equilibrium selection","category_id":21,"category":{"id":21,"name":"economics","display_name":"Economics","description":"Open economic theory statements, published research agendas, and proposed scoped specifications; question provenance and historical priority are qualified per record.","slug":"economics","order_index":21,"created_at":"2026-10-08T00:00:00Z"},"status":"research_question","external_url":null,"legacy_ids":[],"published":false,"statement":"Identify how initial experience and feedback change the probabilities that heterogeneous human learners approach different equilibria, and validate a portable equilibrium-selection model.","economics":{"original_id":52,"field":"Game theory and strategic interaction","kind":"empirical","formalization_basis":"adapted_research_question","source_status":"research_agenda","priority_status":"earliest_found","priority_note":"This is the earliest source in which the relevant agenda was verified during this audit; absolute priority has not been established.","earliest_verified_reference":{"authors":["Drew Fudenberg","David K. Levine"],"title":"Whither Game Theory? Towards a Theory of Learning in Games","year":2016,"url":"https://economics.mit.edu/sites/default/files/2022-09/Whither%20Game%20Theory%20Towards%20a%20Theory%20of%20Learning%20in%20Games.pdf","doi":"10.1257/jep.30.4.151","locator":"Introduction, p.153; conclusion, pp.165–166","evidence_summary":"The essay explicitly calls for more knowledge about learning speed, equilibrium prediction, and tractable models that allow reevaluation."},"current_reference":{"authors":["Drew Fudenberg","David K. Levine"],"title":"Whither Game Theory? Towards a Theory of Learning in Games","year":2016,"url":"https://economics.mit.edu/sites/default/files/2022-09/Whither%20Game%20Theory%20Towards%20a%20Theory%20of%20Learning%20in%20Games.pdf","doi":"10.1257/jep.30.4.151","locator":"Introduction, p.153; conclusion, pp.165–166","evidence_summary":"The essay explicitly calls for more knowledge about learning speed, equilibrium prediction, and tractable models that allow reevaluation."},"formalization_note":"The selection estimand and heterogeneity study are adapted from the essay’s broader predictive agenda, not a named universal selection conjecture."},"statusQualification":"Published question or research agenda verified at the cited locator. The mathematical specification is adapted for this catalogue and includes additional assumptions; source publication does not establish first-ever priority or certify this exact model remains unsolved. The selection estimand and heterogeneity study are adapted from the essay’s broader predictive agenda, not a named universal selection conjecture.","statusReviewedAt":"2026-10-08"}
% FIRST_POSTED: 2026-10-08
% ENTRY_KIND: statement_only
% !TeX program = lualatex
\documentclass[11pt]{article}
\usepackage{fontspec}
\usepackage[a4paper,margin=25mm]{geometry}
\usepackage{amsmath,amssymb,mathtools}
\usepackage{xurl}
\usepackage[hidelinks]{hyperref}
\newcommand{\sref}[2]{\hyperlink{source:#1:#2}{[S#2]}}
\setlength{\emergencystretch}{3em}
\begin{document}
% problemId:30007563
\section*{Learning-based equilibrium selection}\label{30007563}
\subsection{Definitions and mathematical statement}
Consider a preregistered class of finite coordination games \(G\) with a known set \(E_G\) of Nash equilibria. Groups of human players repeatedly interact for \(T\) periods under randomly assigned feedback and initial-history treatment \(Z\). Let \(\widehat\sigma_T\) be the product of empirical action distributions, and fix disjoint neighborhoods \(U_e\) around isolated equilibria \(e\in E_G\). Define the selection probabilities
\[
 q_e(G,Z,P,T)=\Pr_P(\widehat\sigma_T\in U_e\mid do(Z)),
\]
where \(P\) is the sampling population; preserve a residual category for paths lying outside every neighborhood. The target is to estimate these probabilities and identify a portable model relating them to learning-rule heterogeneity, payoff risk, information, and initial experience. Separate a treatment's effect on equilibrium selection from its effect on the speed of approaching any equilibrium by measuring several prespecified horizons. Elicit individual beliefs or learning parameters with measurement error explicitly modeled, and test the selection rule on new games and populations. Random assignment identifies treatment effects within sampled groups; transporting results requires stated population-overlap assumptions. Existing selection results for particular learning dynamics remain valid background. The open agenda concerns predictive success for heterogeneous human learners across contexts, rather than an assertion that equilibrium selection is unknown in every mathematical model.

\par\textbf{Formulation and scope.} The selection estimand and heterogeneity study are adapted from the essay’s broader predictive agenda, not a named universal selection conjecture.

\par\textbf{Open-question provenance.} An explicit question or research agenda was verified at the source locator. This does not by itself certify that every later version remains unresolved \sref{30007563}{1}.

\par\textbf{Historical priority.} This is the earliest source in which the relevant agenda was verified during this audit; absolute priority has not been established.

\subsection{Short English statement}
Identify how initial experience and feedback change the probabilities that heterogeneous human learners approach different equilibria, and validate a portable equilibrium-selection model.

\subsection{Sources}
\begin{itemize}\raggedright
\item[\textnormal{[S1]}]\hypertarget{source:30007563:1}{} Drew Fudenberg, David K. Levine. 2016. Whither Game Theory? Towards a Theory of Learning in Games. Introduction, p.153; conclusion, pp.165–166.
\url{https://economics.mit.edu/sites/default/files/2022-09/Whither%20Game%20Theory%20Towards%20a%20Theory%20of%20Learning%20in%20Games.pdf}.
DOI: 10.1257/jep.30.4.151.
{\small\textit{Earliest verified question/agenda reference; historical priority qualified below. The essay explicitly calls for more knowledge about learning speed, equilibrium prediction, and tractable models that allow reevaluation.}}
\end{itemize}
\end{document}
